\documentclass{article} % TMLR
\usepackage[accepted]{tmlr} % switch to [preprint] before arXiv, [] for submission
\usepackage{amsmath,amssymb,amsfonts}
\usepackage{booktabs}
\usepackage{multirow}
\usepackage{graphicx}
\usepackage{xcolor}
\usepackage[capitalize,noabbrev]{cleveref}
\newcommand{\todo}[1]{\textcolor{red}{[TODO: #1]}}
\newcommand{\best}[1]{\textbf{#1}}
\newcommand{\ours}{\textsc{ProtA}} % protected adapter shorthand

\title{Where to Protect: Insertion Depth Determines the Fate of\\
Importance-Protected Adapters in Continual Learning}

\author{\name Batukaan \"Ozen \email ozenbatukaan@gmail.com \\
\addr Independent Researcher}

\def\month{XX}
\def\year{2026}
\usepackage{url}\def\openreview{\url{https://openreview.net/forum?id=XXXX}}

\begin{document}
\maketitle

\begin{abstract}
Parameter-importance regularization (EWC, MAS) applied to a small adapter on a
frozen pre-trained encoder has recently emerged as a strong recipe for
continual learning (CL). What has not been studied is \emph{where} such a
protected adapter should read the encoder: every published variant taps the
final representation. We present a systematic benchmark of importance-protected
adapters across \emph{insertion depths} of two heterogeneous vision encoders
(supervised ResNet-50, CLIP ViT-B/32), five class-incremental streams
(Split CIFAR-10/100 at 5--20 tasks, Split MNIST, and a heterogeneous
5-dataset stream), five classic CL methods (EWC, MAS, SI, LwF, ER) and two
modern frozen-feature baselines (NCM, RanPAC), all under a single trainer,
validation-based tuning, and five seeds with corrected significance tests.
Three findings emerge. \textbf{(1) Placement is method-agnostic and can improve
the state of the art}: concatenating mid- and late-depth features
(``multi-tap'') lifts not only protected adapters but RanPAC itself, from
82.7\% to 87.3\% on the 5-dataset stream ($+4.6$) and from 95.9\% to 96.2\% on
20-task CIFAR-100 --- a placement-only change to a published method.
\textbf{(2) The optimal placement is encoder-dependent}: late-pair taps win on
ResNet-50, whereas CLIP's final block dominates and early blocks actively hurt.
\textbf{(3) Protection is most critical at early depths}: without it,
sequential fine-tuning collapses by up to 35 points of forgetting at the
earliest taps of both encoders, while MAS restores near-upper-bound accuracy at
every depth. Capacity-matched controls (wider MLP, LoRA) confirm the gains stem
from placement and protection, not parameters, and a class-incremental
evaluation shows the placement effect \emph{amplifies} under the harder
protocol ($+7.7$ for RanPAC vs.\ $+1.0$ task-incremental). Finally, unfreezing
only the least-important 5\% of encoder parameters helps precisely on the most
shifted stream and hurts on near-distribution ones, giving a single actionable
rule: spend both tap depth and encoder plasticity where the distribution shift
is, and nowhere else.
\end{abstract}

\section{Introduction}
Continual learning (CL) on top of frozen pre-trained encoders has become the
dominant regime for vision benchmarks: features are computed once, and a light
trainable module absorbs the task stream
\citep{zhou2023revisiting,mcdonnell2023ranpac,wang2022l2p}. Within this regime,
a particularly simple recipe --- a shared multi-layer adapter protected by a
parameter-importance penalty \citep{kirkpatrick2017ewc,aljundi2018mas} --- has
been rediscovered several times in quick succession
\citep{ewclora2026,ekpc2025,onlinelora2025}. All of these works, however,
attach the adapter to the \emph{final} encoder representation, implicitly
assuming that the last layer is the right place to continually learn.

This paper asks the question none of them ask: \emph{where} should a protected
adapter be inserted? The answer, it turns out, is neither ``at the end'' nor
universal, and getting it right is worth more than switching between
regularizers.

\paragraph{Contributions.}
\begin{itemize}
\item \textbf{A placement map for protected adapters.} We benchmark
  importance-protected adapters at every stage of ResNet-50 and at four depths
  of CLIP ViT-B/32, on five streams, under a unified trainer with
  validation-tuned hyperparameters and 5-seed statistics
  (\cref{sec:setup,sec:depthmap}).
\item \textbf{Placement transfers to the state of the art.} The multi-tap
  configuration discovered on protected adapters also improves RanPAC
  \citep{mcdonnell2023ranpac} on every stream we test --- most notably
  $+4.6$ points on the heterogeneous 5-dataset stream (\cref{sec:transfer}).
\item \textbf{Encoder-dependence and early-depth fragility.} The optimal tap
  set differs qualitatively between a supervised CNN and a contrastive ViT,
  and unprotected sequential training collapses catastrophically precisely at
  the depths where features are most reusable (\cref{sec:depthmap}).
\item \textbf{Controls that keep the comparison honest.} Capacity-matched
  wider adapters and a parameter-matched LoRA variant isolate placement and
  protection from parameter count; we report where protected adapters trail
  RanPAC and show the gap is one of classifier capacity, not forgetting
  (\cref{sec:capacity,sec:honest}).
\item \textbf{A first step past frozen features.} We add \emph{selective
  encoder plasticity} --- unfreezing only the least-important 5\% of encoder
  parameters, protected and anchored --- and show it raises accuracy exactly
  where distribution shift is largest (5-datasets) and harms it where
  pre-training already fits (CIFAR), obeying the same shift rule as placement
  (\cref{sec:beyond}).
\end{itemize}

\section{Related Work}
\label{sec:related}
\paragraph{Importance-based regularization.} EWC \citep{kirkpatrick2017ewc},
SI \citep{zenke2017si} and MAS \citep{aljundi2018mas} penalize movement of
parameters deemed important for previous tasks; LwF \citep{li2016lwf} distills
previous outputs, and experience replay \citep{chaudhry2019er} stores
exemplars. We use all five as same-trainer baselines.

\paragraph{Adapters on frozen backbones.} Residual adapters
\citep{rebuffi2017adapters,houlsby2019adapters} and LoRA \citep{hu2021lora}
transfer frozen models parameter-efficiently. In CL, ADAM/APER
\citep{zhou2023revisiting}, SimpleCIL prototypes, and prompt-based methods
\citep{wang2022l2p,wang2022dualprompt,smith2023coda} adapt frozen ViTs.
Closest to us, EWC-LoRA \citep{ewclora2026}, EKPC \citep{ekpc2025} and
Online-LoRA \citep{onlinelora2025} combine importance penalties with adapters
--- always at the final representation. Our contribution is orthogonal: the
\emph{placement axis}, plus capacity-matched controls none of these report.

\paragraph{Frozen-feature classifiers.} RanPAC \citep{mcdonnell2023ranpac}
projects frozen features through a fixed random expansion and fits a streaming
ridge classifier; LayUP \citep{layup2024} concatenates the last $k$ ViT blocks
in this family, foreshadowing our multi-tap finding within a single encoder
type. NCM \citep{rebuffi2017icarl,zhou2023revisiting} is the prototype
baseline. We run both in our exact setting and \emph{improve} RanPAC with our
placement.

\paragraph{Where-to-adapt \& depth analyses.} Head2Toe \citep{evci2022head2toe}
selects intermediate features for transfer; \citet{ramasesh2021anatomy} and
\citet{davari2022probing} analyze which depths drift under CL;
\citet{lee2023wacv} study pre-trained CL empirically. None map insertion depth
for \emph{protected adapters} across heterogeneous encoders.

\paragraph{Parameter isolation.} PackNet \citep{mallya2018packnet}, CLNP
\citep{golkar2019clnp}, HAT \citep{serra2018hat}, WSN \citep{kang2022wsn},
NISPA \citep{gurbuz2022nispa} and SPU \citep{zhang2024spu} allocate or select
parameters per task. Our selective-plasticity variant (\cref{sec:beyond})
inverts the SPU intuition --- it trains the \emph{least} important encoder
parameters rather than protecting the most important ones --- and unlike
PackNet/WSN stores no per-task masks, relying instead on importance-driven
rotation.

\paragraph{Streams.} The heterogeneous 5-dataset stream follows
\citet{ebrahimi2020adversarial}; we substitute KMNIST for notMNIST for full
torchvision reproducibility (documented in \cref{sec:setup}).

\section{Method and Benchmark Design}
\label{sec:method}
\subsection{Protected adapters at arbitrary insertion depths}
Let $E$ be a frozen encoder with stages $E_1,\dots,E_L$. For a tap set
$\mathcal{T}=\{l_1,\dots,l_m\}$ we pool each tapped stage
($z_l=\mathrm{GAP}(E_l(x))$ for CNNs; the CLS token for ViT blocks), z-score
each tap with train statistics, and concatenate:
$z_\mathcal{T}(x)=[\hat z_{l_1};\dots;\hat z_{l_m}]$.
A single shared 2-layer ReLU MLP adapter $g_\theta$ (hidden width 256, no
normalization layers) maps $z_\mathcal{T}$ to a representation read by
per-task linear heads $h_t$ (task-incremental protocol).

Training task $t$ minimizes
\begin{equation}
\mathcal{L}_t=\mathrm{CE}\big(h_t(g_\theta(z_\mathcal{T}(x))),y\big)
+\alpha\sum_{i} S_i\,(\theta_i-\theta_i^{t-1})^2,
\label{eq:penalty}
\end{equation}
where $S$ is the accumulated importance of the chosen signal --- MAS
($S_i\propto\mathbb{E}\,|\partial\|h(g_\theta(z))\|_2^2/\partial\theta_i|$),
EWC (diagonal Fisher), or SI (path integral) --- normalized by a single global
maximum to preserve cross-tensor ratios. LwF, ER and plain sequential training
(\emph{naive}) share the identical trainer; a pooled multi-head \emph{joint}
run provides the upper bound.

\subsection{Controls}
\textbf{Capacity:} widening the adapter to match multi-tap input
dimensionality (h354) and beyond (h1024), and a parameter-matched linear LoRA
adapter, separate placement effects from parameter count.
\textbf{SOTA baselines:} NCM and RanPAC (random projection dim
$10^4$, streaming Gram + ridge, $\lambda$ tuned on a train split) run on the
\emph{same cached features}, including our multi-tap configurations.

\subsection{Selective encoder plasticity}
\label{sec:method-sp}
To probe beyond frozen features we optionally unfreeze the encoder's
least-important $q$ fraction of parameters (by MAS importance accumulated
jointly through encoder and adapter), hard-masking all gradients outside that
set, protecting the freed subset with its own penalty, and anchoring the
tapped features to the previous task's encoder via an MSE term. The mask is
recomputed each task, so as a parameter's importance rises with use it rotates
out of the trainable set --- capacity rotation without stored masks. We report
this at $q{=}0.05$; \cref{sec:beyond} gives results and ablations.

\section{Experimental Setup}
\label{sec:setup}
\textbf{Streams.} Split CIFAR-10 (5 tasks), Split CIFAR-100 (10 and 20 tasks),
Split MNIST (pixels, sanity), and the heterogeneous 5-dataset stream
CIFAR-10$\to$MNIST$\to$SVHN$\to$Fashion$\to$KMNIST (one dataset per task).
Task-incremental evaluation (task identity at test); metrics: average final
accuracy (AvgAcc), forgetting, plasticity, BWT.
\textbf{Encoders.} ResNet-50 (ImageNet-V2 weights; taps layer1--4:
256/512/1024/2048-d) and CLIP ViT-B/32 (blocks 3/6/9 + final embedding).
Features are cached once; every CL run trains only the adapter and heads.
\textbf{Protocol.} Hyperparameters ($\alpha$, $\lambda_{\mathrm{LwF}}$)
selected on a 10\% train-carved validation split with sweep seed 7, disjoint
from the five reporting seeds $\{42,123,456,789,1337\}$; Welch $t$-tests with
Bonferroni correction across method pairs. Code, tuned configurations, and all
per-seed JSONs are released.

\section{Results}
\subsection{The depth map}
\label{sec:depthmap}
\Cref{tab:depthmap,tab:clipdepth,fig:depthmap} map the fate of an adapter at
every insertion depth. Three regularities hold across both encoders. First,
\textbf{unprotected sequential training fails hardest where it is cheapest}:
naive forgetting climbs monotonically toward early taps (ResNet: 19.4 at
layer4 $\to$ 31.8 at layer1; CLIP: 5.8 at the final embedding $\to$ 35.5 at
block3). Early features are the most \emph{shared} across tasks --- exactly
why overwriting them is catastrophic. Second, \textbf{importance protection
converts every depth into a viable insertion point}, recovering 87--99\% of
the depth-matched joint upper bound (e.g., layer3: 87.2 of 90.7; CLIP block3:
$+25.8$ over naive). The gap between MAS and naive is enormous and
unambiguous ($+16.1$ points on C100 layer3+4, Welch $p<10^{-4}$), whereas the
choice of protection \emph{signal} barely registers: MAS vs.\ EWC is
statistically indistinguishable at every core cell tested (C100 layer3+4
$p{=}0.42$; layer4 $p{=}0.46$). What matters is that \emph{some} importance
protection is present, not which one. Third, \textbf{the placement optimum is
encoder-dependent} (\cref{fig:heatmap}): on ResNet-50 the late pair (layer3+4)
gives the best adapter accuracy (90.2 vs.\ 89.6 at layer4), though on the
near-saturated C100 adapter this particular margin is within seed noise
($p{=}0.10$); the placement effect becomes large and significant in the harder
regimes and, decisively, when transferred to RanPAC (\cref{sec:transfer},
$p{=}2{\times}10^{-4}$). On CLIP the final embedding alone is optimal and every
early-block admixture hurts (final 94.4 vs.\ block3+final 93.3). A supervised
CNN distributes reusable, linearly-separable structure across its last two
stages; a contrastively trained ViT concentrates it in the projected
embedding. ``Where to protect'' therefore has no universal answer --- it must
be searched, and our map is the first such search for protected adapters.
\begin{table}[t]
\centering\small
\caption{Split CIFAR-100 (10 tasks, ResNet-50), AvgAcc\,\% $\pm$ std over 5
seeds (Forgetting\,\% in parentheses). Joint = upper bound. Protection
restores near-upper-bound accuracy at every depth; the collapse of naive
training grows toward early taps.}
\label{tab:depthmap}
\begin{tabular}{lcccc}
\toprule
Tap & Naive & EWC & MAS & Joint \\
\midrule
layer1 & 45.3$\pm$1.0 (31.8) & 64.3$\pm$0.5 (1.5) & 64.9$\pm$0.7 (2.1) & 73.6 \\
layer2 & 55.9$\pm$0.4 (30.6) & 74.1$\pm$0.3 (1.8) & 74.6$\pm$0.4 (1.8) & 81.6 \\
layer3 & 72.1$\pm$1.4 (23.2) & 87.1$\pm$0.4 (2.7) & 87.2$\pm$0.3 (1.5) & 90.7 \\
layer4 & 76.0$\pm$1.9 (19.4) & 89.8$\pm$0.3 (2.1) & 89.6$\pm$0.3 (1.5) & 91.2 \\
layer3+4 & 74.1$\pm$1.5 (22.5) & 89.8$\pm$0.8 (3.4) & \best{90.2$\pm$0.5} (1.8) & 92.2 \\
\bottomrule
\end{tabular}
\end{table}

\begin{table}[t]
\centering\small
\caption{CLIP ViT-B/32, Split CIFAR-100 (10 tasks), 5 seeds: the placement
optimum is encoder-dependent --- the final embedding dominates and adding
early blocks hurts, the opposite of ResNet-50.}
\label{tab:clipdepth}
\begin{tabular}{lccc}
\toprule
Tap & Naive & MAS & EWC \\
\midrule
block3 & 42.0$\pm$1.2 (35.5) & 67.8$\pm$0.4 (1.9) & 58.3$\pm$2.3 (15.4) \\
block9 & 74.4$\pm$2.1 (20.7) & 87.3$\pm$0.5 (1.0) & 87.5$\pm$1.5 (4.1) \\
final & 90.2$\pm$1.1 (5.8) & 94.4$\pm$0.3 (0.4) & \best{95.5$\pm$0.2} (1.0) \\
block9+final & 86.7$\pm$1.0 (10.2) & 93.9$\pm$0.1 (1.0) & 94.6$\pm$0.3 (1.6) \\
block3+final & --- & 93.3$\pm$0.3 (0.7) & --- \\
\bottomrule
\end{tabular}
\end{table}

\begin{figure}[t]
\centering
\includegraphics[width=\linewidth]{figs/fig1_depthmap.pdf}
\caption{AvgAcc vs.\ insertion depth on Split CIFAR-100 (10 tasks). Protection
(MAS/EWC) restores near-upper-bound accuracy at every depth; the naive
collapse grows toward early taps on both encoders, but the placement optimum
differs: late-pair multi-tap for ResNet-50, the final embedding for CLIP.}
\label{fig:depthmap}
\end{figure}

\begin{figure}[t]
\centering
\includegraphics[width=.55\linewidth]{figs/fig2_heatmap.pdf}
\caption{Placement heatmap: MAS-adapter AvgAcc over all ResNet-50 tap pairs
(diagonal = single taps). The late pair L3+4 is the global optimum.}
\label{fig:heatmap}
\end{figure}


\subsection{Placement transfers to RanPAC}
\label{sec:transfer}
If placement were an artifact of our adapter, it would not survive a change of
method family. It does (\cref{tab:transfer}): feeding RanPAC's random
projection from our multi-tap features instead of the final representation
improves it on every stream where the data departs from the encoder's
pre-training distribution, with the gain scaling in the size of that shift ---
$+0.1$/$+0.3$ on near-distribution CIFAR, $+4.6$ (ResNet) and $+1.2$ (CLIP) on
the heterogeneous 5-dataset stream. The 87.3\% (ResNet) and 89.6\% (CLIP)
5-dataset results are, to our knowledge, the strongest reported for this
stream under frozen encoders, and they required \emph{no change to RanPAC
itself}. Placement is a property of the representation, not of the continual
learner attached to it.
\begin{table}[t]
\centering\small
\caption{Applying our multi-tap placement to RanPAC improves it wherever the
stream departs from the pre-training distribution --- a placement-only change
to a published method. The gain scales with domain shift: negligible on
near-ceiling CIFAR, decisive on the heterogeneous 5-dataset stream. 5 seeds
except $^*$.}
\label{tab:transfer}
\begin{tabular}{llccc}
\toprule
Encoder & Stream & RanPAC (last tap) & RanPAC (multi-tap) & $\Delta$ \\
\midrule
ResNet-50 & Split CIFAR-10 & 98.20$\pm$0.06 & \best{98.50$\pm$0.06} & +0.30 \\
ResNet-50 & C100 (10 t.) & 92.11$\pm$0.24 & \best{93.12$\pm$0.13} & +1.01 \\
ResNet-50 & C100 (20 t.) & 95.88$\pm$0.10 & \best{96.15$\pm$0.26} & +0.27 \\
ResNet-50 & 5-Datasets & 82.69$\pm$0.04 & \best{87.28$\pm$0.10} & \best{+4.59} \\
\midrule
CLIP ViT-B/32 & Split CIFAR-10 & 99.54$\pm$0.04 & 99.50$\pm$0.05 & $-$0.04 \\
CLIP ViT-B/32 & C100 (10 t.) & 96.61$\pm$0.06 & 96.40$\pm$0.02 & $-$0.21 \\
CLIP ViT-B/32 & 5-Datasets & 88.42$\pm$0.03 & \best{89.63$\pm$0.03} & \best{+1.21} \\
\bottomrule
\end{tabular}
\end{table}

\subsection{Long and heterogeneous streams; honest SOTA comparison}
\label{sec:honest}
\Cref{tab:long} reports the two hardest streams. Protected adapters dominate
their own family --- on 20-task CIFAR-100 MAS retains 94.5\% with 1.2\%
forgetting where naive training collapses to 78.9\% --- and reach 99.3\% of
their joint upper bound (95.1\%), i.e., forgetting is essentially solved
\emph{within} the adapter family. RanPAC nevertheless scores higher (96.2\%),
and instructively so: it exceeds the adapter's \emph{joint} bound, meaning its
advantage is one of classifier capacity (a $10^4$-dimensional random expansion
with closed-form ridge) rather than of continual-learning ability. The same
pattern holds on the 5-dataset stream (87.3 vs.\ 79.5). We state this
plainly rather than hiding behind favorable cells: under frozen features,
gradient-trained MLP adapters buy \emph{stability} (F$\approx$1 vs.\
RanPAC's 0.4--1.2 with far fewer parameters at train time), while
random-expansion classifiers buy \emph{capacity}. The two are complementary
--- and both obey the same placement map, which is the point of this paper.
\begin{table}[t]
\centering\small
\caption{20-task CIFAR-100 and the 5-dataset stream (ResNet-50, layer3+4
unless noted). RanPAC's ridge classifier exceeds even our adapter's
\emph{joint} upper bound --- the remaining gap is classifier capacity, not
forgetting.}
\label{tab:long}
\begin{tabular}{lcc}
\toprule
Method & C100 20-task & 5-Datasets \\
\midrule
Naive adapter & 78.9$\pm$1.5 (18.5) & 55.1$\pm$1.5 (43.4) \\
MAS adapter & 94.5$\pm$0.2 (1.2)$^{l4}$ & 79.4$\pm$1.4 (7.3) \\
EWC adapter & 94.5$\pm$0.2 (1.5)$^{l4}$ & 81.8$\pm$1.2 (7.2) \\
Joint (adapter UB) & 95.1$\pm$0.2$^{l4}$ & 89.2$\pm$0.1 \\
NCM & 93.6 & 69.4 \\
RanPAC (+our placement) & \best{96.2$\pm$0.3} & 87.3$\pm$0.1 \\
\midrule
CLIP-final MAS adapter & 94.4$^{t10}$ & 82.8$\pm$0.5$^{n2}$ \\
CLIP RanPAC (+our placement) & 96.6$\pm$0.1$^{t10}$ & \best{89.6$\pm$0.0} \\
\bottomrule
\end{tabular}
\end{table}

\subsection{Capacity-matched controls}
\label{sec:capacity}
On Split CIFAR-10 (layer4), widening the adapter from $h{=}256$ (MAS
98.06$\pm$0.11) to the multi-tap-input-matched $h{=}354$ (98.22$^*$) and
further to $h{=}1024$ (97.94$\pm$0.11; EWC 97.42$^*$) does \emph{not}
reproduce the multi-tap gain --- the widest adapters are in fact \emph{worse}
than the narrowest. The same holds on CIFAR-100 at five seeds: $h{=}354$ at
layer4 reaches 89.70$\pm$0.80, below multi-tap's 90.17$\pm$0.51 at matched
parameters. The decisive capacity-independent evidence is \cref{tab:transfer}:
RanPAC's projection dimension is fixed at $10^4$ regardless of input width, so
its multi-tap gains ($+0.3$ to $+4.6$) cannot be a parameter-count effect at
all.

\paragraph{An unexpected control result: identity-initialized residual
adapters.} Our parameter-matched LoRA variant ($h = x + BAx$, $B$
zero-initialized, MAS-protected, rank 144) behaves like the MLP on CIFAR-10
(97.89$\pm$0.06) but \emph{outperforms every gradient-trained MLP
configuration on CIFAR-100} (93.25$\pm$0.12, forgetting 0.10, vs.\ 90.17 for
multi-tap MLP), approaching multi-tap RanPAC (93.12$\pm$0.13). Two caveats
temper the comparison: the residual stream keeps the full 2048-d
representation, so per-task heads are $8\times$ wider than the MLP's, and the
identity initialization means each task starts from an unforgotten linear
probe. We flag this as a finding for the adapter-design axis rather than the
placement axis --- protection composes with \emph{any} of these architectures,
and the placement map applies to all of them.

\subsection{Beyond frozen features: selective encoder plasticity}
\label{sec:beyond}
Frozen features cap every method above; can a \emph{minimally} plastic encoder
raise the ceiling without reintroducing forgetting? We unfreeze only the
bottom 5\% of encoder parameters by accumulated MAS importance (hard gradient
mask, recomputed per task so claimed parameters rotate out), protect that
subset with its own quadratic penalty, and anchor the tapped features to the
previous task's encoder with an MSE term.

The outcome is a clean confirmation of the shift rule from
\cref{sec:transfer}. On the heterogeneous 5-dataset stream, selective
plasticity lifts the MAS adapter from 78.1 to \textbf{81.2} (seed 42, 10
epochs), with the gain coming from \emph{retention}: MNIST end-of-stream
accuracy rises from 83.2 to 93.5 and SVHN from 50.2 to 56.3, while the encoder
moves by only 0.1--0.2\% in relative parameter norm and consecutive trainable
sets overlap 64--87\% (rotation is real but gradual). On near-distribution
streams the same mechanism \emph{costs} accuracy (C100 10-task: 89.1 vs.\
90.2 frozen; 20-task: 93.4 vs.\ 94.5): when pre-training already covers the
stream, any encoder motion is pure risk. Encoder plasticity, like tap depth,
should be spent \emph{where the shift is} --- and nowhere else.

Exposure matters as much as placement: at 20 epochs per task the same
mechanism over-drifts (77.0, below the frozen 78.1), i.e., plasticity has a
\emph{budget} --- brief exposure adapts, prolonged exposure erodes. We
therefore treat epochs-per-task as a first-class hyperparameter of encoder
plasticity, and isolate its contribution with a frozen-encoder control at
matched exposure. \todo{frozen-e10 control, aggressive-adaptation (q=0.2)
and no-anchor ablations running; RanPAC-head combination queued; 5-seed
promotion pending}

\subsection{Task-agnostic (class-IL) evaluation}
\label{sec:classil}
All results above use task-aware evaluation. Merging every head and taking a
global argmax with no task identity at test (\cref{tab:classil}) sharpens two
conclusions. First, \textbf{regularized adapters are task-IL instruments}: MAS
still all but eliminates \emph{forgetting} (6.5\% vs.\ naive's 63.5\% on
C100), but the merged head now suffers heavy cross-task confusion (37.1\%
accuracy), the well-documented failure of regularization-based methods under
class-IL \citep{vandeven2022three}. Second, and more important for this paper,
\textbf{the placement finding not only survives the harder protocol but
amplifies}: giving RanPAC our multi-tap features raises its class-IL accuracy
by $+7.7$ points on C100 (69.1\,$\to$\,76.7), versus $+1.0$ under task-IL ---
the wider the evaluation demands on the representation, the more where-you-tap
matters. RanPAC dominates class-IL outright (76.7 vs.\ the adapter's 37.1),
consistent with \cref{sec:honest}: its advantage is classifier capacity, and
class-IL is where capacity is decisive.

\begin{table}[t]
\centering\small
\caption{Class-incremental (task-agnostic, global-argmax) evaluation,
ResNet-50 layer3+4. Regularized adapters suppress forgetting but not cross-task
confusion; RanPAC dominates, and our multi-tap placement helps it \emph{more}
under class-IL than task-IL ($+7.7$ vs.\ $+1.0$ on C100). $^*$single seed,
5-seed running.}
\label{tab:classil}
\begin{tabular}{llcccc}
\toprule
Stream & & Naive & MAS & NCM & RanPAC \\
\midrule
\multirow{2}{*}{C100 (10 t.)}
 & last tap & 30.8$^*$ (59.9) & 33.9 (4.4) & 63.4 & 69.1$^*$ \\
 & multi-tap & 29.6$\pm$2.2 (63.5) & 37.1$\pm$1.8 (6.5) & 61.9 & \best{76.7$\pm$0.1} \\
\midrule
CIFAR-10 & multi-tap & 62.6$\pm$3.4 (31.2) & 61.7$\pm$2.0 (1.5) & --- & \best{92.4$^*$} \\
5-Datasets & multi-tap & 28.4$^*$ (64.0) & 38.9$^*$ (10.0) & --- & \best{87.0$^*$} \\
\bottomrule
\end{tabular}
\end{table}

\section{Discussion and Limitations}
\label{sec:limits}
\paragraph{What the map buys a practitioner.} Given a new encoder and stream,
our results reduce the design space to one cheap sweep: cache per-stage pooled
features, probe each tap (and the late pair, for CNNs) with any protected
adapter or with RanPAC, and pick by validation. The sweep costs minutes on
cached features --- orders of magnitude less than a single fine-tuning run ---
and, per \cref{tab:transfer}, its outcome transfers across method families.

\paragraph{Why multi-tap helps CNNs but not CLIP.} ResNet's layer3 retains
mid-level structure (texture, parts) that layer4's class-aligned pooling
discards; on streams whose classes differ from ImageNet's, this complements
the final features. CLIP's training objective already compresses transferable
structure into the projected embedding, so early CLS tokens add mostly noise
--- except under strong domain shift (5-Datasets), where even CLIP benefits
from a second late tap. This yields a testable rule of thumb: \emph{the wider
the distribution shift, the deeper into the network the useful taps reach.}

\paragraph{Limitations.} (i) Our primary protocol is task-incremental;
\cref{sec:classil} adds merged-head class-IL results, but stronger class-IL
machinery (task-id inference, calibrated heads) is future work. (ii) Two
encoders; DINOv2-scale
self-supervised models remain to be mapped. (iii) Classic small-image streams
under large pre-trained backbones --- near-distribution CIFAR results sit at
ceiling and we interpret them only through forgetting. (iv) Placement is
searched over stage-level taps, not arbitrary layers. (v) End-of-stream
hyperparameter tuning on a validation split mirrors common practice but is an
oracle relative to strictly-online CL.

\section{Conclusion}
Importance-protected adapters are a strong, simple recipe for continual
learning on frozen encoders --- but \emph{where} they read the encoder has
been an unexamined default. Mapping insertion depth systematically, we find
that placement (i) changes outcomes by more than the choice of regularizer,
(ii) transfers across method families, improving RanPAC by up to $+4.6$ points
with no algorithmic change, (iii) depends qualitatively on the encoder, and
(iv) interacts with distribution shift by a simple rule: the wider the shift,
the deeper the useful taps. We release the full benchmark --- one trainer,
seven methods, two encoders, five streams, five seeds --- so that future
adapter-based CL methods can report \emph{where} they attach, not only
\emph{how} they protect.

\bibliography{refs}
\bibliographystyle{tmlr}

\end{document}
